% VI39-FINITE-MORPHISM-DIVISOR-LAYER
\section{Finite dominant morphisms: controlled pullback and pushforward}\label{sec:vi39-finite-morphism-divisors}
The warning above concerns arbitrary morphisms.  For a finite dominant morphism between normal Noetherian integral schemes, codimension-one behavior is controlled by finite extensions of DVRs, and both pullback and pushforward admit clean formulas.

\begin{proposition}[Zero divisor versus global unit]\label{prop:vi39-zero-divisor-unit}
Let \(X\) be normal Noetherian integral and \(u\in K(X)^\times\).  Then
\[
 \operatorname{div}_X(u)=0
 \quad\Longleftrightarrow\quad
 u\in\Gamma(X,\mathcal O_X)^\times.
\]
\end{proposition}

\begin{proof}
A global unit has order zero at every codimension-one point.  Conversely, suppose every order of \(u\) is zero.  On a normal affine open \(V=\operatorname{Spec}A\), the Noetherian normal domain \(A\) is a Krull domain and
\[
 A=\bigcap_{\operatorname{ht}\mathfrak p=1}A_{\mathfrak p}
 \subset\operatorname{Frac}(A).
\]
Both \(u\) and \(u^{-1}\) have nonnegative order at every height-one prime, hence both lie in every \(A_{\mathfrak p}\), and therefore in \(A\).  Thus \(u\) is a unit on every affine open and hence globally.
\end{proof}

Let \(g:Y\to X\) now be finite dominant, with function-field extension \(L=K(Y)\supset K=K(X)\).  If \(D\subset X\) is a prime divisor and \(E\subset Y\) is a prime divisor lying over \(D\), then the inclusion of DVRs
\[
 \mathcal O_{X,\eta_D}\subset\mathcal O_{Y,\eta_E}
\]
defines the \emph{ramification index}
\[
 e(E/D):=\operatorname{ord}_E(\pi_D),
\]
where \(\pi_D\) is any uniformizer at \(D\).

\begin{definition}[Finite pullback and pushforward]\label{def:vi39-finite-pull-push}
For a prime divisor \(D\) on \(X\), define
\[
 g^*[D]:=\sum_{E\mid D}e(E/D)[E].
\]
For a prime divisor \(E\) on \(Y\), let \(D=g(E)\); finiteness implies that \(D\) is again codimension one.  Define
\[
 g_*[E]:=[\kappa(E):\kappa(D)]\,[D].
\]
Extend both maps linearly to Weil divisors.
\end{definition}

\begin{proposition}[Principal divisors and finite pullback]\label{prop:vi39-finite-pullback-principal}
For every \(f\in K(X)^\times\),
\[
 \operatorname{div}_Y(g^*f)=g^*\operatorname{div}_X(f).
\]
\end{proposition}

\begin{proof}
If \(E\) lies over \(D\), extension of normalized DVR valuations gives
\[
 \operatorname{ord}_E(f)=e(E/D)\operatorname{ord}_D(f).
\]
This is exactly the coefficient of \([E]\) on the right-hand side.
\end{proof}

\begin{proposition}[Principal divisors and the field norm]\label{prop:vi39-finite-pushforward-norm}
For \(h\in L^\times\),
\[
 g_*\operatorname{div}_Y(h)
 =\operatorname{div}_X\!\bigl(N_{L/K}(h)\bigr).
\]
\end{proposition}

\begin{proof}
For every prime divisor \(D\) of \(X\), the valuation formula for the norm in a finite extension of discretely valued fields is
\[
 \operatorname{ord}_D\!\bigl(N_{L/K}(h)\bigr)
 =\sum_{E\mid D}[\kappa(E):\kappa(D)]\operatorname{ord}_E(h).
\]
The right-hand side is precisely the coefficient of \([D]\) in \(g_*\operatorname{div}_Y(h)\).
\end{proof}

\begin{example}[The power map on the affine line]\label{ex:vi39-power-map-affine}
Let
\[
 g:\mathbb A^1_u\longrightarrow\mathbb A^1_x,
 \qquad x=u^m,
\]
with \(m\ge1\).  The divisor \([0]_x\) has one prime divisor above it, namely \([0]_u\), and
\[
 e(0/0)=\operatorname{ord}_{u=0}(u^m)=m.
\]
Hence
\[
 g^*[0]_x=m[0]_u,
\]
which agrees with
\[
 \operatorname{div}(g^*x)=\operatorname{div}(u^m)=m[0]_u.
\]
Also \(g_*[0]_u=[0]_x\) when the residue fields are both \(k\).
\end{example}

\begin{example}[The power map on \(\mathbb P^1\)]\label{ex:vi39-power-map-projective}
The morphism
\[
 [s:t]\longmapsto[s^m:t^m]
\]
is finite.  It is totally ramified at \(0\) and \(\infty\), so
\[
 g^*[0]=m[0],\qquad g^*[\infty]=m[\infty].
\]
Since \(\operatorname{div}(t/s)=[0]-[\infty]\) in the chosen coordinate convention, pulling back gives
\[
 \operatorname{div}\!\left((t/s)^m\right)
 =m[0]-m[\infty]
 =g^*\operatorname{div}(t/s).
\]
This is the finite-morphism version of principal-divisor functoriality.
\end{example}

\begin{warning}[Scope of the finite formulas]\label{warn:vi39-finite-scope}
The formulas in this section use finiteness in an essential way.  They should not be read as defining a pullback of arbitrary Weil divisors for every morphism.  Flat pullback, proper pushforward of cycles, and the full intersection-theoretic formalism belong to later theory.
\end{warning}
